\medskip

%On considère les fonctions $f$ et $g$ suivantes :
%
%\[\begin{array}{l l}
%f \,: &x \longmapsto 4x + 3\\
%g \,:&x \longmapsto 6x
%\end{array}\]
%
%Leurs représentations graphiques $(d_1)$ et $(d_2)$ sont tracées ci-dessous :

\begin{enumerate}
\item %Parmi ces deux fonctions, laquelle représente une situation de proportionnalité ?
C'est l'application linéaire $x \longmapsto g(x) = 6x$ qui représente une situation de proportionnalité.
\item %Calculer l'image de 0 par la fonction $g$.
L'image de 0 par $g$ est $g(0) = 6 \times 0 = 0$ : la droite $(d_1)$ contient l'origine.
\item L'antécédent de $0$ par la fonction $f$ est le nombre $x$ tel que $f(x) = 0$, soit $4x + 3 = 0$ ou $4x = -3$ et $x = -3 \times \dfrac14 = - \dfrac34$.
%Déterminer l’antécédent de $0$ par la fonction $f$.

\begin{center}
\psset{xunit=1.7cm,yunit=0.325cm,arrowsize=2pt 3,comma}
\begin{pspicture*}(-1.5,-1.9)(5.5,20)
\multido{\n=-1.5+0.5}{16}{\psline[linewidth=0.25pt](\n,-1.9)(\n,20)}
\multido{\n=-0+2}{11}{\psline[linewidth=0.25pt](-1.5,\n)(5.5,\n)}
\psaxes[linewidth=1.25pt,Dx=0.5,Dy=2,labelFontSize=\scriptstyle]{->}(0,0)(-1.5,-1.9)(5.5,20)
\psplot[plotpoints=800,linewidth=1.25pt,linecolor=red]{-1.5}{4.5}{x 4 mul 3 add}
\psplot[plotpoints=800,linewidth=1.25pt,linecolor=blue]{-1.5}{4.5}{x 6 mul}
\uput[ul](3,18){\blue $(d_1)$}\uput[dr](3.75,18){\red $(d_2)$}
\end{pspicture*}
\end{center}

\item %Associer à chaque droite la fonction qu'elle représente. Justifier la réponse.
La droite $d_1$représente l'application linéaire donc la droit $d_2$ représente l'application affine $f$.
\item %Déterminer graphiquement les coordonnées du point d’intersection des droites $(d_1)$ et $(d_2)$.
Graphiquement on lit les coordonnées du point commun aux deux droites : (1,5~;~9).

Ceci signifie que $f(1,5) = 4 \times 1,5 +3 = 6 + 3 = 9$ et que $g(1,5) = 6 \times 1,5 = 9$ : les images par $f$ et par $g$ de 1,5 dnt égales à 9.
\end{enumerate}

\bigskip
